\chapter*{Abstract}
\thispagestyle{fancy}
\addcontentsline{toc}{chapter}{Abstract}
\begingroup
\setlength{\emergencystretch}{2em}

This report presents the analytical design and CST full-wave simulation of an
inset-fed rectangular microstrip patch antenna operating near 915~MHz. The antenna
uses a Rogers RT/duroid 5880 substrate with relative permittivity
$\varepsilon_r=2.20$, loss tangent $\tan\delta=0.0009$, and thickness
$h=1.575$~mm. Standard transmission-line-model relations give an initial patch width
of 129.51~mm, effective dielectric constant of 2.1605, fringing extension of
0.834~mm, and physical patch length of 109.79~mm. The documented 50-$\Omega$
feed-line width is 4.89~mm, and the cosine-squared input-resistance estimate gives an
inset depth of 38.13~mm for the final patch length.

The final CST geometry uses $W=129.5$~mm, $L=108.2$~mm, a 4.89~mm feed line,
a 6.89~mm inset-notch width, and a 38.2~mm inset depth. The preserved CST marker
places the matching minimum at 0.9186~GHz with
$S_{11,\min}=-22.14379$~dB. This corresponds to a positive return loss of
22.14379~dB and a frequency deviation of 3.6~MHz, or 0.393\%, from the 915~MHz
target.

At 0.915~GHz, the CST far-field plots display a peak directivity of 7.01~dBi and a
broadside main beam. The $\phi=0^\circ$ cut shows a main-beam direction of
1$^\circ$, a 3-dB angular width of 82.3$^\circ$, and a displayed sidelobe level of
$-19.2$~dB. The $\phi=90^\circ$ cut shows a main-beam direction of 0$^\circ$ and
a 3-dB angular width of 90.8$^\circ$. The study connects closed-form synthesis,
parameterized model construction, resonance adjustment, impedance matching, field
visualization, and radiation verification.

All performance results are analytical or simulated. No fabricated prototype or
measured antenna data are claimed. The absence of raw CST exports, native model files,
and formal mesh- and boundary-convergence records limits the conclusions to the
preserved evidence package.
\par\endgroup
